% 第 6 题程序框图：按固定源内容流矢量层精确值复刻（1 pt → 0.03515 cm，与全库同制）。
% 原点取框图竖轴 x=216.1406 pt、底缘（结束框底）y（顶向下）=688.7070 pt；y 向上。
% 源实测（figure-inventory §1）：竖轴严格 x=216.1406；7 框（开始=圆角 27.988×16.832、
% k=1 矩形 32.035×14.285、k=k+1 矩形 50.660×15.160、k>4? 菱形 76.008×18.879、
% S=sin(kπ/6) 矩形 57.219×28.586（内嵌真分式故双倍高）、输出 S 平行四边形（斜偏 6.133）、
% 结束=圆角 27.988×17.008）；连线 6 竖直（x=0）+「否」支 3 段回折（自菱形右顶点右行、
% 上行、左行，箭尖 (0,162.1992) 恰落于 k=1 框与 k=k+1 框之间竖线的中点——回流到 k=k+1 之前）；
% 全图无虚线、除 7 支实心箭头外无填充；框内文字 STSong/CM 系 10.5 pt，本图用 \small 同制。
\begin{tikzpicture}[x=0.03515cm,y=0.03515cm,>=stealth,every node/.style={font=\small,inner sep=0pt}]
% ===== 7 框（外轮廓照源实测值）=====
\draw[line width=0.398bp,rounded corners=0.1399cm] (-13.9922,198.3359) rectangle (13.9961,215.1679);  % 1 开始
\draw[line width=0.398bp] (-16.0156,169.4843) rectangle (16.0196,183.7695);                            % 2 k=1
\draw[line width=0.398bp] (-25.3281,139.7578) rectangle (25.3321,154.9179);                            % 3 k=k+1
\draw[line width=0.398bp] (0,125.1093)--(38.0039,115.6679)--(0,106.2304)--(-38.0039,115.6679)--cycle;   % 4 菱形 k>4?
\draw[line width=0.398bp] (-28.6093,62.9961) rectangle (28.6094,91.5820);                              % 5 S=sin(kπ/6)
\draw[line width=0.398bp] (-25.3984,31.5742)--(-19.2656,48.4297)--(25.3945,48.4297)--(19.2617,31.5742)--cycle; % 6 输出 S
\draw[line width=0.398bp,rounded corners=0.1399cm] (-13.9922,0) rectangle (13.9961,17.0078);            % 7 结束
% ===== 连线（6 竖直 + 1 回折；箭尖照源落在下一框边线上）=====
\draw[->,line width=0.398bp] (0,198.1367)--(0,183.9687);   % 开始→k=1
\draw[->,line width=0.398bp] (0,169.2851)--(0,155.1172);   % k=1→k=k+1（回折靶线）
\draw[->,line width=0.398bp] (0,139.5586)--(0,125.3906);   % k=k+1→菱形顶
\draw[->,line width=0.398bp] (0,105.9492)--(0,91.7812);    % 菱形底→S=sin（是支）
\draw[->,line width=0.398bp] (0,62.7968)--(0,48.6289);     % S=sin→输出 S
\draw[->,line width=0.398bp] (0,31.3750)--(0,17.2070);     % 输出 S→结束
\draw[->,line width=0.398bp] (38.2852,115.6679)--(52.4531,115.6679)--(52.4531,162.1992)--(0,162.1992); % 否支回折（回流 k=k+1 入线中点）
% ===== 框内文字与分支标签（照源字形中心定位）=====
\node at (0,207.2570) {开始};
\node at (0,175.5170) {\(k=1\)};
\node at (0,146.6540) {\(k=k+1\)};
\node at (0,114.6860) {\(k>4?\)};
\node at (0,76.7790) {\(S=\sin\dfrac{k\pi}{6}\)};
\node at (0,40.4740) {输出 \(S\)};
\node at (0,8.9590) {结束};
\node at (-8.8666,99.4340) {是};
\node at (38.3814,124.5270) {否};
\end{tikzpicture}
